\section{Distributional Forecaster with Conformal Calibration}
\label{sec:method}
\subsection{Encoder and Attention}
Fig.~\ref{fig:arch} shows the forecaster. Each input row is concatenated with a learned embedding of the sine and cosine of the hour of day and the day of week, and a two-layer GRU encoder with hidden size $d=128$~\cite{cho2014learning} maps the window to the states $\mathbf{e}_{t-L+1},\dots,\mathbf{e}_t$. At decoder step $h$ the additive attention~\cite{bahdanau2015neural} forms the context $\mathbf{c}_h=\sum_{j}a_{h,j}\mathbf{e}_j$ with $a_{h,j}\propto\exp\big(\mathbf{v}^\top\tanh(\mathbf{W}_q\mathbf{s}_{h-1}+\mathbf{W}_k\mathbf{e}_j)\big)$, where $\mathbf{s}_{h-1}$ is the top decoder state. The encoder and the attention are those of the baseline. The comparison therefore isolates the decoder and the output.

\begin{figure}[t]
  \centering
  \resizebox{\columnwidth}{!}{% final forecaster: GRU encoder, additive attention, residual autoregressive decoder.
\begin{tikzpicture}[
  font=\footnotesize, >=Latex,
  b/.style={draw, rounded corners=2pt, minimum height=7mm, align=center, fill=#1},
  b/.default=gray!10,
  lnk/.style={->, thick}
]
\node[b=orange!20, text width=24mm] (in) at (0,0) {input window\\$\mathbf{X}_t\in\mathbb{R}^{L\times F}$};
\node[b=yellow!20, text width=24mm] (te) at (0,-1.3) {periodic time\\embedding (hour, day)};
\node[b=blue!12, text width=22mm] (enc) at (3.2,-0.65) {2-layer GRU\\encoder ($d{=}128$)};
\node[b=blue!12, text width=20mm] (att) at (6.2,0.6) {additive\\attention};
\node[b=green!15, text width=22mm] (dec) at (6.2,-0.9) {2-layer GRU\\decoder cell};
\node[b=green!15, text width=20mm] (head) at (6.2,-2.3) {per-target\\heads $\Delta\hat{y}$};
\node[draw, circle, inner sep=1pt, fill=white] (sum) at (3.2,-2.3) {$+$};
\node[b=red!12, text width=24mm] (last) at (0,-2.9) {last observation\\$\mathbf{y}_t$ (step $h{=}1$)};
\node[b=violet!12, text width=22mm] (out) at (3.2,-3.6) {$\hat{\mathbf{y}}_{t+h}$, $h{=}1,\dots,H$};
\draw[lnk] (in) -- (enc);
\draw[lnk] (te) -- (enc);
\draw[lnk] (enc) |- node[pos=0.75, above]{\scriptsize keys} (att);
\draw[lnk] (enc) -- node[above]{\scriptsize $\mathbf{h}_L$} (dec);
\draw[lnk] (att) -- node[right]{\scriptsize context} (dec);
\draw[lnk] (dec) -- (head);
\draw[lnk] (head) -- (sum);
\draw[lnk] (last) -| (sum);
\draw[lnk] (sum) -- (out);
\draw[lnk, dashed] (out.east) -| ++(4.0,0) |- node[pos=0.25, right, align=left]{\scriptsize fed back\\[-1pt]\scriptsize (teacher\\[-1pt]\scriptsize forcing in\\[-1pt]\scriptsize training)} (dec.east);
\draw[lnk, dashed] (dec.north east) to[out=30, in=-30, looseness=1.4] node[right]{\scriptsize query} (att.south east);
\end{tikzpicture}
}
  \caption{Proposed forecaster. The decoder is anchored on the last observation, and each head outputs a mean and five non-crossing quantiles of the change.}
  \label{fig:arch}
\end{figure}

\subsection{Anchored Distributional Decoder}
A two-layer GRU decoder cell takes $\mathbf{c}_h$ and an embedding of the previous mean forecast $\hat{\boldsymbol{\mu}}_{t+h-1}$, with $\hat{\boldsymbol{\mu}}_{t}=\mathbf{y}_t$ read from the input window, and produces $\mathbf{s}_h$. For each target $k$ a head $g_k$ (a two-layer perceptron) outputs six numbers $(m,r,a_1,a_2,b_1,b_2)$, from which
\begin{align}
\hat{\mu}_{t+h,k}&=\hat{\mu}_{t+h-1,k}+m, \qquad \hat{q}_{0.5}=\hat{\mu}_{t+h-1,k}+r,\label{eq:anchor}\\
\hat{q}_{0.25}&=\hat{q}_{0.5}-\varsigma(a_1),\quad \hat{q}_{0.025}=\hat{q}_{0.25}-\varsigma(a_2),\nonumber\\
\hat{q}_{0.75}&=\hat{q}_{0.5}+\varsigma(b_1),\quad \hat{q}_{0.975}=\hat{q}_{0.75}+\varsigma(b_2),\nonumber
\end{align}
with $\varsigma(z)=\log(1+e^{z})+\epsilon$ and $\epsilon=10^{-4}$. The anchoring in~\eqref{eq:anchor} makes persistence the forecast obtained when all head outputs are zero. The network only learns departures from the last value. Because $\varsigma>0$, the five quantiles are strictly ordered for every input, which removes the quantile crossing that separate quantile heads can produce. The decoder feeds back the mean, which keeps the recursion identical to the baseline.

\subsection{Training Loss}
With targets normalized by training statistics, the loss over a batch of $N$ windows is
\begin{equation}
\mathcal{L}=\frac{1}{NHK}\sum_{n,h,k}\Big[(\hat{\mu}_{nhk}-y_{nhk})^2+\frac{1}{5}\sum_{\tau\in\mathcal{T}}\rho_\tau(y_{nhk}-\hat{q}_{\tau,nhk})\Big],
\label{eq:loss}
\end{equation}
with $\mathcal{T}=\{0.025,0.25,0.5,0.75,0.975\}$. By Proposition~\ref{prop:optimal}, the population minimizer of~\eqref{eq:loss} over unrestricted functions of the window is the conditional mean for $\hat{\mu}$ and the conditional $\tau$-quantile for each $\hat{q}_\tau$. These targets are compatible with the ordering constraint, since conditional quantiles are ordered. The point forecast reported for squared-error metrics is $\hat{\mu}$ and for absolute-error metrics $\hat{q}_{0.5}$. Teacher forcing is applied per sample with a probability that decays from 0.5 to 0 over half of the epochs~\cite{bengio2015scheduled}.

\subsection{Conformal Calibration Across Sessions}
The learned quantiles are fitted on training sessions and may under-cover or over-cover new sessions. Let $\mathcal{C}$ be the windows of the calibration sessions and, for each step $h$ and target $k$, define the CQR score~\cite{romano2019conformalized}
\begin{equation}
s_i=\max\{\hat{q}_{\alpha/2}(\mathbf{X}_i)-y_i,\;y_i-\hat{q}_{1-\alpha/2}(\mathbf{X}_i)\},\qquad i\in\mathcal{C}.
\end{equation}
With $c_{h,k}$ the $\lceil(|\mathcal{C}|+1)(1-\alpha)\rceil$-th smallest score, the interval is $[\hat{q}_{\alpha/2}-c_{h,k},\,\hat{q}_{1-\alpha/2}+c_{h,k}]$. If the calibration and test windows are exchangeable, its coverage is at least $1-\alpha$~\cite{romano2019conformalized}. Windows of one session overlap and sessions differ in time of day and load. Exchangeability therefore holds only approximately~\cite{barber2023conformal}, and we report the coverage measured on test sessions. For a model without quantiles we use the split conformal interval $\hat{y}\pm c$ with $c$ the same order statistic of $|y_i-\hat{y}_i|$~\cite{lei2018distribution}.

\subsection{Seed Ensembles and Complexity}
Three models trained with different seeds are combined by averaging their mean and each quantile, which keeps the ordering, and the ensemble is calibrated in the same way~\cite{lakshminarayanan2017simple}. The forecaster has \ParamsDist{} parameters against \ParamsPrev{} for the scalar baseline on OneWeb DL. A forecast needs one encoder pass and $H$ decoder steps, as in the baseline. Algorithm~\ref{alg:cqr} summarizes training, calibration, and inference.

\begin{algorithm}[t]
\caption{Training, calibration, and online use}
\label{alg:cqr}
\begin{algorithmic}[1]
\REQUIRE training, calibration, and test sessions, $\alpha$, seeds $\mathcal{S}$
\FOR{each seed in $\mathcal{S}$}
  \STATE train the forecaster with~\eqref{eq:loss} on training sessions, select the checkpoint with the lowest smoothed calibration-session loss
\ENDFOR
\STATE average the means and quantiles of the seeds
\FOR{each step $h$ and target $k$}
  \STATE compute the CQR scores on calibration windows and the offset $c_{h,k}$
\ENDFOR
\STATE \textbf{online:} for a new window, output $\hat{\mu}$, $\hat{q}_{0.5}$, and $[\hat{q}_{\alpha/2}-c_{h,k},\hat{q}_{1-\alpha/2}+c_{h,k}]$ for all $h,k$
\end{algorithmic}
\end{algorithm}
